% Generated draft: requires mathematical and source review.
\documentclass[11pt]{article}
\usepackage[margin=25mm]{geometry}
\usepackage{fontspec}
\setmainfont{Latin Modern Roman}
\usepackage{amsmath,amssymb,mathtools,unicode-math}
\setmathfont{Latin Modern Math}
\usepackage{enumitem,needspace,xurl,hyperref}
\setlength{\parindent}{0pt}
\setlength{\parskip}{5pt}
\setlength{\emergencystretch}{3em}
\newcommand{\N}{\mathbb{N}}
\newcommand{\Z}{\mathbb{Z}}
\newcommand{\Q}{\mathbb{Q}}
\newcommand{\R}{\mathbb{R}}
\newcommand{\C}{\mathbb{C}}
\newcommand{\F}{\mathbb{F}}
\newcommand{\A}{\mathbb{A}}
\newcommand{\PP}{\mathbb{P}}
\newcommand{\E}{\mathbb{E}}
\newcommand{\eps}{\varepsilon}
\newcommand{\dd}{\,\mathrm{d}}
\begin{document}
\section*{UnsolvedMath ID: 30006004}
\textbf{Source problem code:} catalogue entry
\subsection{Definitions and mathematical statement}
Fix a prime \(p\), an integer \(n\geq 1\), and integers \(h_1\geq\cdots\geq h_n\). Put
\[
G_p=\operatorname{Gal}(\overline{\Q}_p/\Q_p),\qquad I_p\subseteq G_p
\]
for the absolute Galois group and its inertia subgroup. Let \(\omega:G_p\to\overline{\F}_p^{\times}\) be the reduction modulo \(p\) of the cyclotomic character; its restriction to inertia has order \(p-1\).

For a continuous crystalline representation \(\rho:G_p\to\operatorname{GL}_n(E)\), where \(E/\Q_p\) is finite, choose a \(G_p\)-stable lattice \(T\) over the ring of integers \(\mathcal O_E\). If \(k_E\) is its residue field, define its semisimplified reduction by
\[
\overline\rho=\left((T/\mathfrak m_E T)\otimes_{k_E}\overline{\F}_p\right)^{\mathrm{ss}}.
\]
Here semisimplification means the direct sum of the Jordan--H\"older constituents. Its isomorphism class is independent of the stable lattice. Coefficient extensions are allowed when considering representations with values in \(\overline{\Q}_p\). The source does not specify the sign convention for Hodge--Tate weights; the weights below must be interpreted in the convention of the cited report.

The general classification problem is to determine, up to isomorphism, the set
\[
\mathcal R_p(h_1,\ldots,h_n)=
\left\{[\overline\rho]:\rho\text{ is crystalline with Hodge--Tate multiset }\{h_1,\ldots,h_n\}\right\}.
\]
There is no restriction to a short interval of Hodge--Tate weights in this general question.

The supplied excerpt reports the following necessary conditions in the case where the semisimplified residual representation is a direct sum of characters. In that case write
\[
\overline\rho|_{I_p}\cong\bigoplus_{i=1}^n\omega^{m_i}|_{I_p}.
\]
The exponents initially determine residue classes modulo \(p-1\). According to the excerpt's Theorem 1, they admit integer representatives \(m_1\geq\cdots\geq m_n\) such that
\[
\sum_{i=1}^k m_i\leq\sum_{i=1}^k h_i\quad(1\leq k\leq n),\qquad
\sum_{i=1}^n m_i=\sum_{i=1}^n h_i,
\]
and
\[
\exists\sigma\in S_n\ \forall i\in\{1,\ldots,n\},\qquad
m_i\equiv h_{\sigma(i)}\pmod p.
\]
The congruences modulo \(p\) concern the chosen integer representatives, whereas the inertia characters depend on their classes modulo \(p-1\).

The particular sharpness question asks whether these conditions are best possible for \(n>2\): does every multiset of inertia characters admitting such integer representatives occur for a crystalline representation with the prescribed Hodge--Tate weights, in the residual character setting intended by the source? The supplied material does not specify whether sharpness concerns only inertia characters or also each prescribed full residual representation, including its unramified character values. These are distinct formulations, and the latter is not asserted here. Nor does this restricted question replace the general classification problem, which also includes residual representations with higher-dimensional irreducible constituents.

For context, the excerpt attributes to Gee--Liu--Savitt [GLS14] the assertion that, in the residual character setting, if every \(h_i\in[0,p]\), the exponents can be chosen, after reordering, equal to the Hodge--Tate weights. It attributes to Berger--Breuil [BB05] a two-dimensional example with Hodge--Tate multiset \(\{0,h\}\), where \(p+1\leq h\leq 2p-1\): the two inertia types are
\[
1\oplus\omega^h\quad\text{and}\quad\omega\oplus\omega^{h-1},
\]
and both are reported to occur for each such \(h\). The representative multisets permitted by Theorem 1 in this example are reported as \(\{h,0\}\) and \(\{p,h-p\}\), consistently with
\[
\omega^p\oplus\omega^{h-p}=\omega\oplus\omega^{h-1}.
\]
The excerpt credits Levin--Wang-Erickson [LWE20] with the partial-sum condition and describes the congruence condition as new in that report. It states that, using the Breuil--M\'ezard conjecture, Theorem 1 can be shown best possible for \(n=2\) for arbitrary Hodge--Tate weights. This qualification is retained; no unconditional sharpness theorem is asserted here.

The transcription loses distinctions between \(\rho\) and \(\overline\rho\), prints \(n=0\) in the two-dimensional example, and writes equality at \(i=n\) in a condition indexed by \(k\). The residual-character interpretation, the two-dimensional reading of the example, and equality at \(k=n\) are contextual interpretations requiring source review. The supplied dated curation labels the problem partially solved, says general classification and sharpness for \(n>2\) remain unresolved, and reports further constraints from 2026 work. These are supplied assessments, not independently verified mathematical or literature conclusions.
\subsection{Short English statement}
For fixed \(p\), dimension, and Hodge--Tate weights, classify the semisimplified reductions modulo \(p\) of crystalline representations of \(G_{\Q_p}\). In the residual character case, determine whether the reported partial-sum inequalities and congruences on integer representatives completely describe the possible inertia characters for \(n>2\). The source's transcription, Hodge--Tate convention, and precise meaning of sharpness require review. Its two-dimensional sharpness discussion uses the Breuil--M\'ezard conjecture; the supplied partially solved status is dated curation.
\subsection{Sources}
\begin{itemize}
\item Local UnsolvedMath registry, supplied complete source record: permanent ID \(30006004\), title \emph{Mod-p Reductions of Crystalline Galois Representations}, source problem code \texttt{catalogue entry}, in the Oberwolfach Reports Open Problems collection. The registry supplies the labels \texttt{partially\_solved} and \texttt{corrected\_verified}; these labels are not independent mathematical verification.
\item Supplied report citation: \emph{Arithmetic Geometry (2025)}, Oberwolfach Reports, DOI \texttt{10.4171/owr/2024/33}, pp.\ 1855--1912. \url{https://doi.org/10.4171/owr/2024/33}. The excerpt identifies the contribution as \emph{The Reduction modulo p of Crystalline Breuil--Kisin Modules}, Toby Gee, joint work with Mark Kisin. No individual proposer is identified in the registry. Its proposed year and year tag are 2024, whereas the supplied report title says 2025; this discrepancy is unresolved.
\item Supplied literature source, labeled ``Gee--Kisin, crystalline Breuil--Kisin reductions'': \url{https://ems.press/content/serial-article-files/50046}.
\item Supplied literature source, labeled ``Bhatt--Gee--Kisin, mod p reductions'': \url{https://arxiv.org/abs/2607.08660}.
\item References named in the excerpt: Gee--Liu--Savitt [GLS14], Berger--Breuil [BB05], and Levin--Wang-Erickson [LWE20]. Full bibliographic entries are not supplied. The Breuil--M\'ezard conjecture is invoked without a specific citation or formulation.
\item Literature assessment date: \texttt{2026-08-22}. This is supplied dated curation, not independently verified. It reports strong constraints for the character case, two-dimensional sharpness, unresolved general classification and higher-dimensional sharpness, and additional constraints from 2026 work. The excerpt's explicit qualification concerning use of the Breuil--M\'ezard conjecture must accompany the two-dimensional claim. No linked source was visited for this draft.
\end{itemize}
\end{document}
